% =====================================================================
% Section: Conclusion  (outline.md §5)
% Draft v1 (2026-06-27). Summarizes contributions and results, notes the
%   transferability of DPR as a generic content-routing module, and future work.
% Consistent with the honest framing of the body (matches/slightly surpasses
%   CATANet with a comparable lightweight cost profile; ablation is now a 250k
%   3-seed study with C1 EMA comparison and C4 variance reduction supported).
% No numeric claims beyond those established in Sec. 3.
% =====================================================================
\section{Conclusion}
\label{sec:conclusion}

We presented Dynamic Prototype Routing (DPR), a drop-in replacement for the
center-based Token Aggregation Block of CATANet that keeps the backbone and tensor
interface unchanged while remaining in a comparable lightweight operating regime. DPR
generates semantic prototypes directly from the current tokens, decouples their
confirmation through a learnable query bank, matches tokens back with a learnable
temperature, and propagates explicit per-token confidence through ordering, gating,
and a soft fallback. A usage-balancing regularizer keeps prototype use spread out. On
five standard benchmarks at $\times2/\times3/\times4$, the resulting DPRNet maintains
CATANet-level reconstruction quality with small gains on most benchmark cells and a
comparable lightweight cost profile. Across three seeds, the ablations attribute the
effect of each design element and show that dynamic prototypes outperform a
re-implemented EMA history center. The routing diagnostics further provide
mechanism-level evidence that the temperature remains unsaturated and that prototype
usage becomes measurably more balanced---and more stable across seeds---in all blocks.
Because DPR preserves the external interface of the routing block, it is not tied to
this backbone and can serve as a confidence-aware content-routing module for other
token-aggregation architectures. Future work includes remeasuring all comparison models
for cross-method runtime under a single implementation protocol, and extending DPR to
larger training sets and real-world or arbitrary-scale degradations.
